\centerline{\bf \Huge Олимпиада III}
\title{Олимпиада III}

1. Компания ЛМШ организовала однокруговой футбольный турнир, где за победу в матче дают 3 балла, за ничью 1 балл, а за поражение пельмени. Всего участвовало 5 команд: команда Анджура, команда Расула, команда Очира, команда Лососей и команда шиншилл.
Победила команда Анджура, и оказалось, что они набрали столько баллов, сколько у всех остальных вместе взятых. Сколько игр закончились в ничью?

2. Расул и Ази праздновали день рождения и решили поделить праздничный торт. Торт был круглый, и по краям торта было расположено 20 свечей. Они договорились, что за раз можно провести разрез между двумя свечками(соединяющий две свечки прямой) так, чтоб не было пересечений с другими разрезами. Тот, кто не сможет сделать разрез, отдает весь торт другому. Кто победит при правильной тактике, если первая ходит Ази?

3. Компания ЛМШ организовала второй однокруговой футбольный турнир с теми же командами, что из первой задачи. Известно что команда Очира в играх против команд Шиншилл и Лососей в сумме забила на 8 голов больше чем в матчах против команд Расула и Анджура вместе взятых. Так же известно, что Лососям команда Очира забивала в три раза чаще, чем Шиншиллам. В конце турнира у команды Очира оказалось строго меньше 12 забитых мячей. Какое максимальное число побед у команды Очира? 

4. Шиншилла и Лосось пришли на кухню и увидели, что остались 5 пиал по 40 бëригов. Они решили, что каждый возьмёт по пиале, а три оставшиеся достанутся победителю интересной игры. Правила простые: за ход можно взять две пиалы и переложить все пельмени из этих пиал в 4 новые пиалы(в каждой пиале должно быть хотя бы по одному бёригу). Проигрывает тот, кто не способен сделать ход. Кто может выиграть при любых ходах противника, если начинает Шиншилла? 

5. Колдун Вячеслав Вуб взял в плен всех детей ЭлМШ и написал на доске 1967 натуральных чисел и сказал, что если они найдут среди этих чисел сколько-нибудь, сумма которых делится на 1967, то он их отпустит, а если нет, то заставит делать ему курсовую работу. Смогут ли дети спастись от злого Вуба?